\documentclass[11pt]{article}
\usepackage[a4paper,margin=25mm]{geometry}
\usepackage{amsmath,amssymb,amsthm,mathtools}
\usepackage[hidelinks]{hyperref}
\allowdisplaybreaks[2]
\setlength{\emergencystretch}{2em}
\title{A vanishing universal-input capacity fraction for additive noise}
\author{Leo Gao\\[-1pt]{\small AI-assisted research}}
\date{November 26, 2026 \quad$\cdot$\quad Research draft}
\begin{document}
\maketitle
\begin{abstract}
We show that no power-constrained input guarantees a positive universal fraction of additive-noise channel capacity. Even when the input is chosen using the noise variance $N$, the worst-case capacity fraction is $O(1/log N)$ as $N$ grows. The noise can be chosen to have a symmetric $C^1$ compactly supported density or a smooth, everywhere-positive density.
\end{abstract}
\section{The universal-input question}

In their 2004 paper, \emph{A Gaussian Input Is Not Too Bad}~\cite{ZE}, Ram Zamir and Uri Erez established that Gaussian inputs lose at most half a bit per sample relative to the capacity of an arbitrary power-constrained additive-noise channel.

In Section VII, equation (24), they explicitly posed a stronger question: does there exist an input \(X\) and an absolute constant \(\alpha>0\) such that

\[
I(X;X+Z)>\alpha C(Z;P)
\]

for every independent noise distribution \(Z\), where

\[
C(Z;P)=
\sup_{\mathbb E[U^2]\le P} I(U;U+Z)?
\]

Twenty years later, Lampros Gavalakis, Ioannis Kontoyiannis and Mokshay Madiman revisited this question in \emph{The Entropic Doubling Constant and Robustness of Gaussian Codebooks for Additive-Noise Channels}~\cite{GKM}.

Their Theorem 8 proved the partial result

\[
I(X_G;X_G+Z)
\geq
\frac{\mathrm{SNR}}{3\mathrm{SNR}+2}C(Z;P),
\]

where \(X_G\sim\mathcal N(0,P)\). The guaranteed fraction tends to zero at low SNR, leaving the original universal-constant question unresolved by their theorem.

\section{Main theorem}

\textbf{Theorem (Universal-input impossibility).} Normalize the transmitter's power constraint to \(\mathbb E[X^2]\leq1\), and define

\[
R(N)=
\sup_{\mathbb E[X^2]\leq1}
\inf_{Z\in\mathcal Z_N}
\frac{I(X;X+Z)}{C(Z;1)},
\]

where \(\mathcal Z_N\) is the class of independent, zero-mean noise distributions with variance \(N\) and finite differential entropy.

Then, as \(N\to\infty\),

\[
\boxed{R(N)=O\left(\frac1{\log N}\right)}
\]

and consequently,

\[
\boxed{\lim_{N\to\infty}R(N)=0.}
\]

The conclusion remains valid when the noise is restricted to symmetric, continuously differentiable, compactly supported probability densities, or alternatively to everywhere-positive, infinitely differentiable densities.

\section{Proof}

All logarithms below are natural. The choice of logarithm base does not affect capacity ratios.

\subsection{Lemma 1 --- A channel with a scale-dependent capacity}

Let \(W\) have probability density

\[
f(w)=
\begin{cases}
\dfrac{15}{16}(1-w^2)^2,& |w|\leq1,\\
0,& |w|>1.
\end{cases}
\]

Direct integration gives

\[
\mathbb E[W]=0,\qquad
\operatorname{Var}(W)=\frac17.
\]

Consider the additive-noise channel with noise \(aW\), where \(a\geq2\).

Choose a binary input

\[
U=
\begin{cases}
3a,&\text{with probability }p=\dfrac1{9a^2},\\
0,&\text{with probability }1-p.
\end{cases}
\]

It satisfies the power constraint exactly:

\[
\mathbb E[U^2]=1.
\]

The two conditional output supports are

\[
[-a,a]\quad\text{and}\quad[2a,4a].
\]

They are disjoint. Consequently, \(U\) can be recovered without error, and

\[
I(U;U+aW)=H_b(p),
\]

where \(H_b\) is binary entropy.

Therefore,

\[
\begin{aligned}
C(aW;1)
&\geq H_b(p)\\
&\geq p\log(1/p)\\
&\geq\frac{2\log a}{9a^2}.
\end{aligned}
\]

We obtain

\[
\boxed{C(aW;1)\geq\frac{2\log a}{9a^2}.}
\]

This lower bound is achieved by an input whose nonzero amplitude depends explicitly on the noise scale. \(\square\)

\subsection{Lemma 2 --- A universal input cannot exploit every scale}

Define the scale-localization function

\[
\omega(r)=
\begin{cases}
r,&0\leq r\leq\frac12,\\
r^{-2},&r>\frac12.
\end{cases}
\]

For every input \(X\) with \(\mathbb E[X^2]\leq1\) and every \(a\geq2\), I claim that

\[
\boxed{
I(X;X+aW)
\leq
\frac{
72+120\log a\,
\mathbb E[X^2\omega(|X|/a)]
}{a^2}.
}
\]

\textbf{Proof.} Introduce an auxiliary reference density

\[
q_a(w)=(1-\delta)f(w)+\delta g(w),
\]

where

\[
\delta=a^{-2},
\qquad
g(w)=\frac1{\pi(1+w^2)}.
\]

The Cauchy component is only a comparison distribution used to bound relative entropy; it is not part of the physical noise channel.

Write \(f_u(y)=f(y-u)\). The information-radius inequality gives

\[
I(X;X+aW)
\leq
\mathbb E\left[D(f_{X/a}\Vert q_a)\right].
\]

This follows directly from the relative-entropy chain rule.

We will establish the pointwise bound, with \(r=|u|\),

\[
D(f_u\Vert q_a)
\leq
2a^{-2}+70r^2+
120\log a\,r^2\omega(r).
\tag{1}
\]

\textbf{Case A: \(0<r\leq1/2\).}

Split the integral defining relative entropy into the central region

\[
A_r=\{w:|w|\leq1-2r\}
\]

and its complement within the support of \(f\).

On \(A_r\),

\[
q_a(w+u)\geq(1-\delta)f(w+u).
\]

Let \(L(w)=\log f(w)\) inside the support. Then

\[
L''(w)=-\frac{4(1+w^2)}{(1-w^2)^2}.
\]

For every point \(t\) between \(w\) and \(w+u\),

\[
|L''(t)|\leq
\frac{32}{(1-|w|)^2}.
\]

Taylor's theorem therefore yields

\[
L(w)-L(w+u)
\leq-uL'(w)+
\frac{16r^2}{(1-|w|)^2}.
\]

The linear term integrates to zero over \(A_r\), since \(f\) is even, \(L'\) is odd and \(A_r\) is symmetric.

Furthermore,

\[
\int_{-1}^{1}
\frac{f(w)}{(1-|w|)^2}\,dw
=\frac{35}{8}.
\]

Using \(-\log(1-\delta)\leq2\delta\), the central contribution to the divergence is at most

\[
2\delta+70r^2.
\]

For the boundary region, observe that

\[
f(w)\leq\frac{15}{4}(1-|w|)^2.
\]

Consequently,

\[
\Pr(|W|>1-2r)\leq20r^3.
\]

On this region, \(|w+u|\leq3/2\), so the reference density satisfies

\[
\frac{f(w)}{q_a(w+u)}
\leq\frac{10}{\delta}=10a^2.
\]

Since \(a\geq2\),

\[
\log(10a^2)\leq6\log a.
\]

Thus the boundary contribution is at most

\[
120r^3\log a.
\]

Combining the two regions proves

\[
D(f_u\Vert q_a)
\leq2a^{-2}+70r^2+120r^3\log a,
\]

which is precisely (1) in Case A. The case \(r=0\) follows directly from \(q_a\geq(1-\delta)f\).

\textbf{Case B: \(r>1/2\).}

For every \(w\) in the support of \(f\),

\[
q_a(w+u)\geq
\frac{\delta}{\pi(1+(w+u)^2)}.
\]

Using \(|w|\leq1\),

\[
\frac{f(w)}{q_a(w+u)}
\leq9a^2(1+r^2).
\]

It follows that

\[
\begin{aligned}
D(f_u\Vert q_a)
&\leq2\log a+\log9+\log(1+r^2)\\
&\leq2\log a+70r^2.
\end{aligned}
\]

Because \(r^2\omega(r)=1\) in this case, inequality (1) follows here as well.

Finally, substitute \(r=|X|/a\) into (1) and take expectations. The power constraint gives

\[
I(X;X+aW)
\leq
\frac{72+
120\log a\,\mathbb E[X^2\omega(|X|/a)]}{a^2}.
\]

This proves the lemma. \(\square\)

\subsection{Lemma 3 --- The multiscale obstruction}

Consider dyadic amplitude scales

\[
a_j=2^j,\qquad j=m,\ldots,2m.
\]

For every real \(x\),

\[
\sum_{j\in\mathbb Z}
\omega\left(\frac{|x|}{2^j}\right)
\leq\frac{19}{3}.
\tag{2}
\]

To see this, divide the terms at the threshold \(r=1/2\).

The terms for \(r\leq1/2\) form a geometric series with total at most \(1\). The remaining terms are inverse squares forming a geometric series with total at most \(16/3\).

Multiplying (2) by \(X^2\), taking expectations and using the power constraint gives

\[
\sum_{j=m}^{2m}
\mathbb E\left[
X^2\omega\left(\frac{|X|}{a_j}\right)
\right]
\leq\frac{19}{3}.
\]

There must therefore exist at least one \(j\in\{m,\ldots,2m\}\) for which

\[
\boxed{
\mathbb E[X^2\omega(|X|/a_j)]
\leq\frac{19}{3(m+1)}.
}
\]

Combining Lemmas 1 and 2 at this particular scale yields

\[
\begin{aligned}
\frac{I(X;X+a_jW)}{C(a_jW;1)}
&\leq
\frac{324}{\log a_j}
+540\mathbb E[X^2\omega(|X|/a_j)]\\
&\leq
\frac{324}{m\log2}
+\frac{3420}{m+1}.
\end{aligned}
\]

The right-hand side tends to zero as \(m\to\infty\).

This already disproves a variance-independent universal input. The next step establishes the stronger result. \(\square\)

\subsection{Lemma 4 --- Concealing the relevant scale from the transmitter}

Fix a noise variance \(N\geq16\). Let

\[
m=\left\lfloor\frac{\log_2N}{4}\right\rfloor.
\]

All the candidate scales in Lemma 3 satisfy

\[
a_j^2\leq N.
\]

Select the scale \(a\) whose existence is guaranteed by that lemma.

Set

\[
v=N-\frac{a^2}{7},
\qquad
L=10(a+\sqrt N).
\]

Independently of \(X\) and \(W\), introduce a random variable \(J\) with distribution

\[
J=
\begin{cases}
-L,&\text{probability }\dfrac{v}{2L^2},\\
0,&\text{probability }1-\dfrac{v}{L^2},\\
L,&\text{probability }\dfrac{v}{2L^2}.
\end{cases}
\]

Define the physical channel noise as

\[
Z=aW+J.
\]

By construction,

\[
\mathbb E[Z]=0,
\]

and

\[
\operatorname{Var}(Z)
=\frac{a^2}{7}+v=N.
\]

The Markov chain

\[
X\longrightarrow X+aW
\longrightarrow X+aW+J
\]

implies, by the data-processing inequality,

\[
I(X;X+Z)\leq I(X;X+aW).
\tag{3}
\]

Crucially, adding \(J\) does not destroy the capacity lower bound from Lemma 1.

Use the same binary input \(U\in\{0,3a\}\). Its conditional output supports now consist of six separated intervals.

For any fixed value of \(J\), the two output intervals remain disjoint. Intervals corresponding to different values of \(J\) are also disjoint because \(L>6a\).

Hence \(U\) remains perfectly recoverable:

\[
C(Z;1)\geq H_b(p)
\geq\frac{2\log a}{9a^2}.
\tag{4}
\]

Combining (3), (4) and the preceding lemmas gives

\[
\boxed{
\frac{I(X;X+Z)}{C(Z;1)}
\leq
\frac{324}{m\log2}
+\frac{3420}{m+1}.
}
\]

This bound holds for every admissible input \(X\), even if its distribution was chosen with full knowledge of \(N\).

The adversarial noise has exactly the prescribed variance, a symmetric continuously differentiable density, compact support and finite differential entropy. Its capacity is finite by the maximum-entropy bound.

Taking the supremum over all admissible inputs establishes

\[
R(N)=O\left(\frac1{\log N}\right).
\]

Therefore,

\[
\boxed{\lim_{N\to\infty}R(N)=0.}
\]

\(\blacksquare\)

\section{A stronger extension: everywhere-positive smooth noise}

One potential objection is that the construction exploits exact gaps in the support of the noise density.

That restriction can be removed.

Let \(G\) be an independent standard Gaussian and replace the constructed noise with

\[
Z_\tau=aW+J+\tau G,
\]

adjusting the variance of \(J\) to

\[
v=N-\frac{a^2}{7}-\tau^2.
\]

For sufficiently small positive \(\tau\), this is a valid construction with exactly the required variance.

The noise now possesses an everywhere-positive, infinitely differentiable density. By data processing,

\[
I(X;X+Z_\tau)\leq I(X;X+aW).
\]

The binary input from Lemma 1 is no longer perfectly recoverable, but its decoding-error probability is bounded by

\[
P_e\leq
\Pr\left(|\tau G|\geq\frac a4\right).
\]

This tends to zero as \(\tau\to0\).

By Fano's inequality, choose \(\tau\) sufficiently small that

\[
I(U;U+Z_\tau)\geq\frac12H_b(p).
\]

The preceding universal-input upper bound consequently changes by at most a factor of two.

\section*{Disclosure}
This manuscript was prepared with AI assistance. The arguments and priority claims are open to independent review.

\begin{thebibliography}{9}
\bibitem{ZE} R. Zamir and U. Erez, A Gaussian input is not too bad, IEEE Transactions on Information Theory 50 (2004), 1362--1367. \href{https://doi.org/10.1109/TIT.2004.828153}{Source.}
\bibitem{GKM} L. Gavalakis, I. Kontoyiannis and M. Madiman, The entropic doubling constant and robustness of Gaussian codebooks for additive-noise channels, IEEE Transactions on Information Theory (2024). \href{https://arxiv.org/abs/2403.07209}{Source.}
\end{thebibliography}
\end{document}
